\question[3]
a. Gib für die beteiligten Verweisboxen die Attributwerte nach Ablauf der untenstehenden Anweisungen an.
(Angabe in der Form:  inhalt = 'e', oben = a, unten = b) \\
b. Welche Ausgabe erscheint auf der Konsole?  
\begin{lstlisting}
class VerweisBox:
    def __init__(self, inhalt, unten=None, oben=None):
        self.inhalt = inhalt
        self.unten = unten
        self.oben = oben
        if unten is not None: self.unten.oben = self
        if oben  is not None: self.oben.unten = self

    def __str__(self):
        return self.inhalt

a = VerweisBox('a')
b = VerweisBox('b')
c = VerweisBox('c',a,b)
d = VerweisBox('d',b,c)
print(a.unten, b.unten.oben, b.oben.oben.oben, c.unten, c.oben.oben.unten.unten.unten)

\end{lstlisting}

\begin{solutionbox}{4.5cm}
\begin{lstlisting}
   Die Werte der Attribute:
   inhalt = 'a', unten = None, oben = c
   inhalt = 'b', unten = c, oben = d
   inhalt = 'c', unten = d, oben = b
   inhalt = 'd', unten = b, oben = c
   
   Die print-Ausgabe:
   None b b d d
\end{lstlisting}
\end{solutionbox}
